% GENERATED FILE. Edit canonical lessons, then run npm run generate:books.
% canonical-source-hash: fnv1a64:22d2b27fb734803e
% canonical-lessons: KA-C197-kavalugara, KA-C197-kuli, KA-C197-kandaktar, KA-C197-mukhyopadhyaya, KA-C197-karyadarsi, KA-R197-first-pass-words-from-chapters-194-195, KA-R197-second-pass-words-from-chapters-196-197

\chapter{Watchman, Daily wages, Bus conductor, Headmaster, Secretary}
\label{ch:ka-watchman-daily-wages-bus-conductor-headmaster-secretary}
\hlchaptermodality{197}

\begin{chapteropening}
\textbf{By the end of this chapter:} \emph{I can say a watchman, daily wages, a bus conductor, a headmaster and a secretary.}

Complete the last lesson of chapter 197: I can say a watchman, daily wages, a bus conductor, a headmaster and a secretary.
\end{chapteropening}

\section[\texorpdfstring{\kn{ಕಾವಲುಗಾರ} (kāvalugāra)}{ಕಾವಲುಗಾರ (kāvalugāra)}]{\kn{ಕಾವಲುಗಾರ} (kāvalugāra) — a watchman, a security guard}

\label{lesson:KA-C197-kavalugara}



\begin{quote}
Before the new one: say the Kannada for a spectator, then the Kannada for a carpenter.
\end{quote}

\subsection*{You'll want to know: \kn{ಕಾವಲುಗಾರ}}
\textbf{\kn{ಕಾವಲುಗಾರ}} — \emph{kāvalugāra} — \textquotedblleft{}a watchman, a security guard\textquotedblright{}.

\begin{grammarlens}[title={What you've built}]
One more word for this chapter.
\end{grammarlens}

\subsection*{Your turn}
\begin{itemize}
\raggedright
  \item \emph{Say it:} \emph{kāvalugāra}
  \item \emph{Say it:} \emph{kāvalugāra}, once more
  \item \emph{Say it:} say \emph{kāvalugāra}
  \item \emph{Recall:} say the Kannada for a spectator, then the Kannada for a carpenter, then say \emph{kāvalugāra} again
\end{itemize}

\begin{tcolorbox}[breakable,colback=teal!4,colframe=teal!35!black,title={Before you move on}]
What does \kn{ಕಾವಲುಗಾರ} mean? (\textquotedblleft{}A watchman, a security guard\textquotedblright{}.) Say it once more.
\end{tcolorbox}
\section[\texorpdfstring{\kn{ಕೂಲಿ} (kūli)}{ಕೂಲಿ (kūli)}]{\kn{ಕೂಲಿ} (kūli) — daily wages; a labourer}

\label{lesson:KA-C197-kuli}



\begin{quote}
Before the new one: say the Kannada for a carpenter, then the Kannada for a watchman.
\end{quote}

\subsection*{You'll want to know: \kn{ಕೂಲಿ}}
\textbf{\kn{ಕೂಲಿ}} — \emph{kūli} — \textquotedblleft{}daily wages; a labourer\textquotedblright{}.

\begin{grammarlens}[title={What you've built}]
One more word for this chapter.
\end{grammarlens}

\subsection*{Your turn}
\begin{itemize}
\raggedright
  \item \emph{Say it:} \emph{kūli}
  \item \emph{Say it:} \emph{kūli}, once more
  \item \emph{Say it:} say \emph{kūli}
  \item \emph{Recall:} say the Kannada for a carpenter, then the Kannada for a watchman, then say \emph{kūli} again
\end{itemize}

\begin{tcolorbox}[breakable,colback=teal!4,colframe=teal!35!black,title={Before you move on}]
What does \kn{ಕೂಲಿ} mean? (\textquotedblleft{}Daily wages; a labourer\textquotedblright{}.) Say it once more.
\end{tcolorbox}
\section[\texorpdfstring{\kn{ಕಂಡಕ್ಟರ್} (kaṇḍakṭar)}{ಕಂಡಕ್ಟರ್ (kaṇḍakṭar)}]{\kn{ಕಂಡಕ್ಟರ್} (kaṇḍakṭar) — a bus conductor}

\label{lesson:KA-C197-kandaktar}



\begin{quote}
Before the new one: say the Kannada for a watchman, then the Kannada for daily wages.
\end{quote}

\subsection*{You'll want to know: \kn{ಕಂಡಕ್ಟರ್}}
\textbf{\kn{ಕಂಡಕ್ಟರ್}} — \emph{kaṇḍakṭar} — \textquotedblleft{}a bus conductor\textquotedblright{}.

\begin{grammarlens}[title={What you've built}]
One more word for this chapter.
\end{grammarlens}

\subsection*{Your turn}
\begin{itemize}
\raggedright
  \item \emph{Say it:} \emph{kaṇḍakṭar}
  \item \emph{Say it:} \emph{kaṇḍakṭar}, once more
  \item \emph{Say it:} say \emph{kaṇḍakṭar}
  \item \emph{Recall:} say the Kannada for a watchman, then the Kannada for daily wages, then say \emph{kaṇḍakṭar} again
\end{itemize}

\begin{tcolorbox}[breakable,colback=teal!4,colframe=teal!35!black,title={Before you move on}]
What does \kn{ಕಂಡಕ್ಟರ್} mean? (\textquotedblleft{}A bus conductor\textquotedblright{}.) Say it once more.
\end{tcolorbox}
\section[\texorpdfstring{\kn{ಮುಖ್ಯೋಪಾಧ್ಯಾಯ} (mukhyōpādhyāya)}{ಮುಖ್ಯೋಪಾಧ್ಯಾಯ (mukhyōpādhyāya)}]{\kn{ಮುಖ್ಯೋಪಾಧ್ಯಾಯ} (mukhyōpādhyāya) — a headmaster}

\label{lesson:KA-C197-mukhyopadhyaya}



\begin{quote}
Before the new one: say the Kannada for daily wages, then the Kannada for a bus conductor.
\end{quote}

\subsection*{You'll want to know: \kn{ಮುಖ್ಯೋಪಾಧ್ಯಾಯ}}
\textbf{\kn{ಮುಖ್ಯೋಪಾಧ್ಯಾಯ}} — \emph{mukhyōpādhyāya} — \textquotedblleft{}a headmaster\textquotedblright{}.

\begin{grammarlens}[title={What you've built}]
One more word for this chapter.
\end{grammarlens}

\subsection*{Your turn}
\begin{itemize}
\raggedright
  \item \emph{Say it:} \emph{mukhyōpādhyāya}
  \item \emph{Say it:} \emph{mukhyōpādhyāya}, once more
  \item \emph{Say it:} say \emph{mukhyōpādhyāya}
  \item \emph{Recall:} say the Kannada for daily wages, then the Kannada for a bus conductor, then say \emph{mukhyōpādhyāya} again
\end{itemize}

\begin{tcolorbox}[breakable,colback=teal!4,colframe=teal!35!black,title={Before you move on}]
What does \kn{ಮುಖ್ಯೋಪಾಧ್ಯಾಯ} mean? (\textquotedblleft{}A headmaster\textquotedblright{}.) Say it once more.
\end{tcolorbox}
\section[\texorpdfstring{\kn{ಕಾರ್ಯದರ್ಶಿ} (kāryadarśi)}{ಕಾರ್ಯದರ್ಶಿ (kāryadarśi)}]{\kn{ಕಾರ್ಯದರ್ಶಿ} (kāryadarśi) — a secretary}

\label{lesson:KA-C197-karyadarsi}



\begin{quote}
Before the new one: say the Kannada for a bus conductor, then the Kannada for a headmaster.
\end{quote}

\subsection*{You'll want to know: \kn{ಕಾರ್ಯದರ್ಶಿ}}
\textbf{\kn{ಕಾರ್ಯದರ್ಶಿ}} — \emph{kāryadarśi} — \textquotedblleft{}a secretary\textquotedblright{}.

\begin{grammarlens}[title={What you've built}]
One more word for this chapter.
\end{grammarlens}

\subsection*{Your turn}
\begin{itemize}
\raggedright
  \item \emph{Say it:} \emph{kāryadarśi}
  \item \emph{Say it:} \emph{kāryadarśi}, once more
  \item \emph{Say it:} say \emph{kāryadarśi}
  \item \emph{Recall:} say the Kannada for a bus conductor, then the Kannada for a headmaster, then say \emph{kāryadarśi} again
\end{itemize}

\begin{tcolorbox}[breakable,colback=teal!4,colframe=teal!35!black,title={Before you move on}]
What does \kn{ಕಾರ್ಯದರ್ಶಿ} mean? (\textquotedblleft{}A secretary\textquotedblright{}.) Say it once more.
\end{tcolorbox}
\section[(dialogue)]{Words from chapters 194-195 — a first pass}

\label{lesson:KA-R197-first-pass-words-from-chapters-194-195}



\begin{quote}
Say the words for a headmaster and a secretary.
\end{quote}

\subsection*{Your turn}
Say these aloud:

\begin{itemize}
\raggedright
  \item a salary, a customer, a tailor, housework, a boss
  \item a fee, an account, interest, tax, a bill
\end{itemize}

\begin{tcolorbox}[breakable,colback=teal!4,colframe=teal!35!black,title={Before you move on}]
A salary? (\textbf{\kn{ಸಂಬಳ}}, \emph{sambaḷa}.) A customer? (\textbf{\kn{ಗ್ರಾಹಕ}}, \emph{grāhaka}.) A tailor? (\textbf{\kn{ದರ್ಜಿ}}, \emph{darji}.) Housework? (\textbf{\kn{ಮನೆಗೆಲಸ}}, \emph{manegelasa}.) A boss? (\textbf{\kn{ಮೇಲಧಿಕಾರಿ}}, \emph{mēladhikāri}.) A fee? (\textbf{\kn{ಶುಲ್ಕ}}, \emph{śulka}.) An account? (\textbf{\kn{ಖಾತೆ}}, \emph{khāte}.) Interest? (\textbf{\kn{ಬಡ್ಡಿ}}, \emph{baḍḍi}.) Tax? (\textbf{\kn{ತೆರಿಗೆ}}, \emph{terige}.) A bill? (\textbf{\kn{ಬಿಲ್}}, \emph{bil}.)
\end{tcolorbox}
\section[(dialogue)]{Words from chapters 196-197 — a second pass}

\label{lesson:KA-R197-second-pass-words-from-chapters-196-197}



\begin{quote}
Say the words for a writer and a director.
\end{quote}

\subsection*{Your turn}
Say these aloud:

\begin{itemize}
\raggedright
  \item a writer, a director, a player, a spectator, a carpenter
  \item a watchman, daily wages, a bus conductor, a headmaster, a secretary
\end{itemize}

\begin{tcolorbox}[breakable,colback=teal!4,colframe=teal!35!black,title={Before you move on}]
A writer? (\textbf{\kn{ಲೇಖಕ}}, \emph{lēkhaka}.) A director? (\textbf{\kn{ನಿರ್ದೇಶಕ}}, \emph{nirdēśaka}.) A player? (\textbf{\kn{ಆಟಗಾರ}}, \emph{āṭagāra}.) A spectator? (\textbf{\kn{ಪ್ರೇಕ್ಷಕ}}, \emph{prēkṣaka}.) A carpenter? (\textbf{\kn{ಬಡಗಿ}}, \emph{baḍagi}.) A watchman? (\textbf{\kn{ಕಾವಲುಗಾರ}}, \emph{kāvalugāra}.) Daily wages? (\textbf{\kn{ಕೂಲಿ}}, \emph{kūli}.) A bus conductor? (\textbf{\kn{ಕಂಡಕ್ಟರ್}}, \emph{kaṇḍakṭar}.) A headmaster? (\textbf{\kn{ಮುಖ್ಯೋಪಾಧ್ಯಾಯ}}, \emph{mukhyōpādhyāya}.) A secretary? (\textbf{\kn{ಕಾರ್ಯದರ್ಶಿ}}, \emph{kāryadarśi}.)
\end{tcolorbox}
